\documentclass[11pt]{scrartcl}
\usepackage{evan}
\addtolength{\textheight}{2ex}
\lhead{AMC 12 Preparation}
\rhead{Berkeley Math Circle}

\begin{document}
\setlength{\droptitle}{-2em}
\title{AMC 12 Preparation}
\maketitle

Here are some assorted problems from some short-answer contests.  The answer choices were added by me.

Problems are in roughly increasing order of difficulty.  The last problem in each section is intended to be very difficult (probably harder than the AMC would permit).

\newcommand{\choices}[5]{\par $\hspace*{1.337ex} \textbf{(A)}\ #1 \hfill \textbf{(B)}\ #2 \hfill \textbf{(C)}\ #3 \hfill \textbf{(D)}\ #4 \hfill \textbf{(E)}\ #5 \hspace*{1.337ex}$}

\section{Problems}
\subsection{Algebra}
\begin{enumerate}[\bfseries {A}1.]
	\ii % ANS: 4
		(NIMO 2011) If the answer to this problem is $x$, then compute the value of $\frac{x^2}{8} + 2$.
		\choices{4}{6}{8}{10}{12}
	\ii % ANS: 45
		(OMO Winter 2013) A permutation $a_1, a_2, \ldots, a_{13}$ of the numbers from $1$ to $13$ is given such that $a_i > 5$ for $i=1,2,3,4,5$.  Determine the maximum possible value~of \[ a_{a_1} + a_{a_2} + a_{a_3} + a_{a_4} + a_{a_5}. \]
		\choices{42}{45}{51}{53}{55}
	\ii % ANS: 178
		(OMO Winter 2013) Determine the absolute value of the sum \[ \lfloor 2013\sin{0^\circ} \rfloor + \lfloor 2013\sin{1^\circ} \rfloor + \cdots + \lfloor 2013\sin{359^\circ} \rfloor, \] where $\lfloor x \rfloor$ denotes the greatest integer less than or equal to $x$.
		\choices{90}{178}{179}{180}{360}
	\ii % ANS: 2013^2
		(NIMO 2013) The infinite geometric series of positive reals $a_1, a_2, \dots$ satisfies
		\[ 1 = \sum_{n=1}^\infty a_n = -\frac{1}{2013} + \sum_{n=1}^{\infty} \text{GM}(a_1, a_2, \dots, a_n) = \frac{1}{N} + a_1 \]
		where $\text{GM}(x_1, x_2, \dots, x_k) = \sqrt[k]{x_1x_2\cdots x_k}$ denotes the geometric mean.  Compute $N$.
		\choices{2011^2}{2012^2}{2013^2}{2014^2}{2015^2}
	\ii % ANS: 1/9
		(NIMO 2012) A number is called \emph{purple} if it can be expressed in the form $\frac{1}{2^a 5^b}$ for positive integers $a > b$. Find the sum of all purple numbers.
		\choices{\frac{1}{13}}{\frac{1}{12}}{\frac{1}{11}}{\frac{1}{10}}{\frac{1}{9}}
	\ii % ANS: 11
		(OMO Winter 2012) Let $a,b,c$ be the roots of the cubic $x^3 + 3x^2 + 5x + 7$. Given that $P$ is a cubic polynomial such that $P(a)=b+c$, $P(b) = c+a$, $P(c) = a+b$, and $P(a+b+c) = -16$, find $P(0)$.
		\choices{9}{11}{13}{16}{21}
	\ii % ANS: 999
		(NIMO 2012) Select reals $x_1, x_2, \cdots, x_{333} \ge -1$, and let $S_k = x_1^k + x_2^k + \cdots + x_{333}^k$.  If $S_2 = 777$, compute the least possible value of $S_3$.
		\choices{499}{800}{999}{1110}{1331}
\end{enumerate}

\subsection{Combinatorics}
\begin{enumerate}[\bfseries {C}1.]
	\ii % ANS: 30
		(OMO Winter 2013) David has a collection of 40 rocks, 30 stones, 20 minerals and 10 gemstones. An operation consists of removing three objects, no two of the same type. What is the maximum number of operations he can possibly perform?
		\choices{0}{10}{20}{30}{40}
	\ii % ANS: 37/192
		(Math Prize 2009) Consider a fair coin and a fair 6-sided die.  The die begins with the number 1 face up.  A \textit{step} starts with a toss of the coin: if the coin comes out heads, we roll the die; otherwise (if the coin comes out tails), we do nothing else in this step.  After 5 such steps, what is the probability that the number 1 is face up on the die?
		\choices{\frac{35}{216}}{\frac{35}{192}}{\frac{1}{6}}{\frac{37}{216}}{\frac{37}{192}}
	\ii % ANS: 57
		(Math Prize 2010) Lynnelle took 10 tests in her math class at Stanford. Her score on each test was an integer from 0 through 100. She noticed that, for every four consecutive tests, her average score on those four tests was at most 47.5. What is the largest possible average score she could have on all 10 tests?
		\choices{47.5}{50}{52.5}{54.5}{57}
	\ii % ANS: 511
		(OMO Winter 2013) A set of 10 distinct integers $S$ is chosen. Let $M$ be the number of nonempty subsets of $S$ whose elements have an even sum. What is the minimum possible value of $M$?
		\choices{255}{256}{511}{512}{1023}
	\ii % ANS: 1866
		(NIMO 2013) Tom has a scientific calculator. Unfortunately, all keys are broken except for one row: \verb$1$, \verb$2$, \verb$3$, \verb$+$ and \verb$-$.
		Tom presses a sequence of $5$ random keystrokes; at each stroke, each key is equally likely to be pressed.  The calculator then evaluates the entire expression, yielding a result of $E$.  Find the expected value of $E$. \par (Note: Negative numbers are permitted, so \verb$13-22$ gives $E = -9$. Any excess operators are parsed as signs, so \verb$-2-+3$ gives $E=-5$ and \verb$-+-31$ gives $E = 31$.  Trailing operators are discarded, so \verb$2++-+$ gives $E=2$.  A string consisting only of operators, such as \verb$-++-+$, gives $E=0$.)
		\choices{1866}{1867}{1868}{1869}{1870}
	\ii % ANS: 870
		(NIMO 2011) For the NEMO, Kevin needs to compute the product \[ 9 \times 99 \times 999 \times \cdots \times 999999999. \]  Kevin takes exactly $ab$ seconds to multiply an $a$-digit integer by a $b$-digit integer. Compute the minimum number of seconds necessary for Kevin to evaluate the expression together by performing eight such multiplications.
		\choices{729}{810}{849}{870}{900}
	\ii % ANS: 113
		(OMO Winter 2013) Beyond the Point of No Return is a large lake containing 2013 islands arranged at the vertices of a regular $2013$-gon.  Adjacent islands are joined with exactly two bridges.  Christine starts on one of the islands with the intention of burning all the bridges.  Each minute, if the island she is on has at least one bridge still joined to it, she randomly selects one such bridge, crosses it, and immediately burns it. Otherwise, she stops. \par If the probability Christine burns all the bridges before she stops can be written as $\frac{m}{n}$ for relatively prime positive integers $m$ and $n$, find the remainder when $m+n$ is divided by $1000$.
		\choices{113}{114}{115}{116}{117}
\end{enumerate}

\subsection{Geometry}
\begin{enumerate}[\bfseries {G}1.]
	\ii % ANS: 270
		(NIMO 2012) Hexagon $ABCDEF$ is inscribed in a circle.  If $\measuredangle ACE = 35^{\circ}$ and $\measuredangle CEA = 55^{\circ}$, then compute the sum of the degree measures of $\angle ABC$ and $\angle EFA$.
		\choices{90}{150}{240}{270}{300}
	\ii % ANS: 75
		(OMO Winter 2013) Let $A$, $B$, and $C$ be distinct points on a line with $AB=AC=1$. Square $ABDE$ and equilateral triangle $ACF$ are drawn on the same side of line $BC$. What is the degree measure of the acute angle formed by lines $EC$ and $BF$?
		\choices{15}{30}{45}{60}{75}
	\ii % ANS: 4sqrt(5)
		(OMO Fall 2012) In triangle $ABC$ let $D$ be the foot of the altitude from $A$. Suppose that $AD = 4$, $BD = 3$, $CD = 2$, and $AB$ is extended past $B$ to a point $E$ such that $BE = 5$. Determine the value of $CE$.
		\choices{8}{4\sqrt{5}}{9}{3\sqrt{11}}{6\sqrt{3}}
	\ii % ANS: 10
		(Math Prize 2011) Let $\triangle ABC$ be a triangle with $AB = 3$, $BC = 4$, and $AC = 5$.  Let $I$ be the center of the circle inscribed in $\triangle ABC$.  What is the product of $AI$, $BI$, and $CI$?
		\choices{4\sqrt{2}}{4\sqrt{3}}{10}{11}{9\sqrt{2}}
	\ii % ANS: 47/14
		(OMO Fall 2012) In trapezoid $ABCD$, $AB < CD$, $AB\perp BC$, $AB\parallel CD$, and the diagonals $AC$, $BD$ are perpendicular at point $P$. There is a point $Q$ on ray $CA$ past $A$ such that $QD\perp DC$. Compute $\frac{BP}{AP} - \frac{AP}{BP}$ given that \[ \frac{QP} {AP}+\frac{AP} {QP} = \left( \frac{51}{14}\right)^4 - 2. \]
		\choices{3}{\frac{47}{14}}{\sqrt{14}}{2\sqrt{7}}{\frac{57}{14}}
	\ii % ANS: 37
		(NIMO 2012) In quadrilateral $ABCD$, $AC = BD$ and $\measuredangle B = 60^\circ$. Denote by $M$ and $N$ the midpoints of $\overline{AB}$ and $\overline{CD}$, respectively. If $MN = 12$ and the area of quadrilateral $ABCD$ is 420, then compute $AC$.
		\choices{36}{37}{38}{39}{40}
	\ii % ANS: 23
		(NIMO 2013) Let $AXYZB$ be a convex pentagon inscribed in a semicircle of diameter $AB$.  Suppose that $AZ-AX=6$, $BX-BZ=9$, and that $AY=12$ and $BY=5$.  Find the greatest integer not exceeding the perimeter of quadrilateral $OXYZ$, where $O$ is the midpoint of $AB$.
		\choices{21}{22}{23}{24}{25}
\end{enumerate}

\subsection{Number Theory}
\begin{enumerate}[\bfseries {N}1.]
	\ii % ANS: 2
		(OMO Winter 2013) The number $123454321$ is written on a blackboard. Evan walks by and erases some (but not all) of the digits, and notices that the resulting number (when spaces are removed) is divisible by $9$. What is the fewest number of digits he could have erased?
		\choices{0}{1}{2}{3}{4}
	\ii % ANS: 2500
		(OMO Fall 2012) Let $\operatorname{lcm} (a,b)$ denote the least common multiple of $a$ and $b$. Find the sum of all positive integers $x$ such that $x\le 100$ and $\operatorname{lcm}(16,x) = 16x$.
		\choices{50}{1000}{2500}{2550}{5050}
	\ii % ANS: 20
		(OMO Fall 2012)  Define a sequence of integers by $T_1 = 2$ and for $n\ge2$, $T_n = 2^{T_{n-1}}$. Find the remainder when $T_1 + T_2 + \cdots + T_{256}$ is divided by 255.
		\choices{0}{8}{12}{20}{24}
	\ii % ANS: 1256, 1257 trap
		(OMO Fall 2012) Find the number of integers $a$ with $1\le a\le 2012$ for which there exist nonnegative integers $x,y,z$ satisfying the equation \[x^2(x^2+2z) - y^2(y^2+2z)=a.\]
		\choices{1255}{1256}{1257}{1258}{1259}
	\ii % ANS: 44
		(NIMO 2012) For how many positive integers $2 \le n \le 500$ is $n!$ divisible by $2^{n-2}$?
		\choices{40}{41}{42}{43}{44}
	\ii % ANS: 5
		(NIMO 2013) For each integer $k\ge2$, the decimal expansions of the numbers $1024,1024^2,\dots,1024^k$ are concatenated, in that order, to obtain a number $X_k$.  (For example, $X_2 = 10241048576$.)  If \[ \frac{X_n}{1024^n} \] is an odd integer, find the smallest possible value of $n$, where $n\ge2$ is an integer.
		\choices{2}{3}{4}{5}{6}
	\ii % ANS: 671
		(OMO 2013) For positive integers $n$, let $s(n)$ denote the sum of the squares of the positive integers less than or equal to $n$ that are relatively prime to $n$. Find the greatest integer less than or equal to \[ \sum_{n\mid 2013} \frac{s(n)}{n^2}, \] where the summation runs over all positive integers $n$ dividing $2013$.
		\choices{668}{669}{670}{671}{672}
\end{enumerate}

\section{Problem Sources}
\begin{itemize}
	\ii OMO: \url{http://www.aops.com/Forum/resources.php?c=182&cid=242}.
	\par The Online Math Open is a team-based online contest where teams of four are given about a week to solve several problems.  There are two tests per year.
	\ii NIMO: \url{http://www.internetolympiad.org/}.
	\par See also \url{http:/www.aops.com/Forum/resources.php?c=182&cid=194}.
	\par The NIMO is a series of online contests.  There is an individual contest roughly every month, in addition to a proof-based contest during the winter.
	\ii Math Prize: \url{http://mathprize.atfoundation.org/}
	\par Math Prize for Girls is an annual contest for high school girls.  It's held at MIT.
	\ii Past AMC problems: \url{http://www.aops.com/Forum/resources.php?c=182}.
\end{itemize}

\end{document}
